\section{Evaluation and Verification Protocol}
\label{app:protocol}

\paragraph{Data partitions and leakage control.}
We distinguish experience used for acquisition and training,
development data used for controller verification, and
held-out data used for final evaluation. The intended unit
of separation is the originating episode: segments from the
same episode should remain in one partition even when they
are associated with different skills. For QA, normalized
questions are deduplicated across partitions before answers
are attached. Question--answer matching provides an additional
check rather than replacing question-level separation.
The final evaluation pool is excluded from gate tuning,
checkpoint selection, and prompt development.

\paragraph{Sequential probe construction.}
In ALFWorld, the consolidator groups identical
observation--action--prompt trajectories across skill labels,
preserves the reserved groups across rounds, and excludes
trajectories already present in replay. Other sequential
environments reserve a deterministic tail of successful
memory records per skill after sorting by node identifier.
Reserved records are excluded from the current consolidation
batch. Observation--action probes include up to 16 recent
history entries and are subject to a probe-step cap.

\paragraph{Verification conditions.}
The candidate arm uses the adapted model without episodic
retrieval; the reference arm uses the base model with ETG.
Both arms use the same task interface and scoring procedure.
The implemented reference arm retrieves from the pre-release
store. Removing a record from the training batch does not
necessarily remove it from this retrieval store, so this
condition is distinct from source-excluded retrieval.

For sequential tasks, probes measure agreement with recorded
actions using the same environment prompt builder and action
parser. Agreement is a local behavioral diagnostic: it can
depend on the action classes represented in the probe and
does not capture all valid alternative actions or subsequent
environment states.

For the QA release analysis in
Figure~\ref{fig:verification-release}, each question is counted
as successful when answer F1 is at least 0.5. The skill-level
score is the fraction of successful development questions.
Release requires this fraction to exceed 50\%, together with
the relative-performance and training-evidence checks.
These answer-level scores are distinct from next-action
agreement.


\paragraph{Statistical unit.}
Paired comparisons match methods by environment world or
task instance. Bootstrap draws resample the same instances
for all compared methods and recompute the reported metric.
For Crafter, each sampled world retains its complete
achievement vector; progress or geometric score is then
computed from the resulting achievement rates.

Paired differences are summarized with percentile 95\%
bootstrap intervals. For evaluations with multiple evolution
seeds, seed-level variation is distinguished from
instance-level uncertainty, with hierarchical resampling
where applicable. Single-checkpoint analyses quantify
variation over evaluation instances rather than variation
across independently trained models. The fraction of
bootstrap draws favoring a method is descriptive and is
not reported as a hypothesis-test $p$-value.

\paragraph{Interaction and compute accounting.}
We distinguish student acquisition, teacher acquisition,
verification, consolidation, and final evaluation costs.
Serving measurements include all model calls and tokens
used by each method, including reflection and repeated
generation. Last-attempt success and best-of-attempt success
are different outcomes and are identified separately when
reported. Matched-data controls reuse the same acquired
experience; this does not imply equal end-to-end acquisition
or training cost.

\paragraph{Decision and memory accounting.}
Release decisions are associated with the evaluated
checkpoint and the supporting training and probe evidence.
The relevant quantities include student mastery, available
teacher demonstrations, realized training examples, probe
size, candidate and reference scores, and the configured
relative threshold and absolute floor.

Memory accounting distinguishes records available for
retrieval, records moved to a non-retrievable archive, and
separately retained training data. Counts before and after
release are reported for the affected stores. Archival
release reduces retrieval availability without necessarily
reducing total storage.

\paragraph{LLM-only control.}
For QA, LLM-only denotes one greedy zero-shot chain-of-thought
generation from the unadapted student using the question and
the benchmark-supplied context, without episodic retrieval,
additional search, or an adapter. Interactive controls with
handcrafted strategy hints or action substitutions are
identified separately rather than labeled LLM-only.

\subsection{Controlled evaluation of memory release}
\label{app:release-controls}

\paragraph{Shared experimental conditions.}
We use the fixed Qwen3-8B SFT checkpoint and the initial
store of 600 HotpotQA experience records analyzed in
Figure~\ref{fig:verification-release}.
All conditions share the checkpoint, skill assignments,
prompt construction, retrieval settings, and decoding
configuration. No additional parameter updates are
performed during this comparison.
Development probes determine release decisions before
held-out evaluation; test results are not used to revise
these decisions or their thresholds.

\paragraph{Release-policy controls.}
We compare the following policies:
\begin{itemize}
    \item \textbf{Retain all.}
    All experience records remain retrievable, and all
    skills use the fixed adapted model with ETG.

    \item \textbf{Immediate release.}
    Skills that contributed valid training examples
    release their eligible records after training,
    without post-training behavioral verification.

    \item \textbf{Verified release.}
    Eligible records are released only for skills that
    satisfy the training-evidence checks, relative
    verification threshold, and absolute performance
    floor.

    \item \textbf{Random release.}
    Release selections are sampled without using probe
    scores and matched to the verified policy's
    retrievable record count.
    We repeat the random selection to measure variation
    due to the choice of released memories.
\end{itemize}

For every policy, skills selected for release use the
same adapted model without episodic retrieval.
Other skills use that model with the remaining ETG.
Adapter activation is fixed rather than independently
recomputed for each policy.
Record eligibility, shared-record handling, and any
configured residual retention are applied consistently
across conditions.
Released records are moved to a non-retrievable archive;
the comparison therefore measures retrieval reduction
rather than total-storage reduction.

\paragraph{Held-out performance and uncertainty.}
We evaluate every condition on the same held-out questions
and report aggregate answer EM and token F1.
For each release policy, we compute the paired F1 difference
relative to retaining all records.
Paired bootstrap draws resample test questions using the
same indices for both conditions and recompute the score
difference, yielding percentile 95\% confidence intervals.

For random release, we report performance averaged over
the sampled release selections and account for both
selection-level and question-level variation through
hierarchical bootstrap resampling.
These intervals quantify uncertainty for the fixed
checkpoint, rather than variation across training seeds.